\section{Background and Problem Setting}
\label{sec:background}

\subsection{Setting and corrections}
\label{sec:background-corrections}

A fine-tuned classifier maps an input $x$ to logits $z(x) \in \mathbb{R}^K$, with $K = 2$ for offensive language and
$K = 3$ for sentiment. It is trained on a source distribution $P$ and deployed on a target distribution $Q$ of another
language variety. We observe labelled source splits and, on the target, unlabelled inputs and possibly $n$ labelled
calibration items; every target quantity is evaluated on a disjoint target test split. Every correction considered here is a
map applied to the logits of a fixed model; what distinguishes them is which part of the decision they can move.

Temperature scaling (TS; \citealp{guo2017calibration}) predicts $\mathrm{softmax}(z/T)$, with $T$ fitted on the source
calibration split (\emph{source TS}) or on target labels. With $n$ target labels we fit $\log T$ by maximum a
posteriori estimation under a Gaussian prior centred on the source estimate (\emph{shrinkage TS}). Bias-corrected
temperature scaling (BCTS; \citealp{alexandari2020maximum}) predicts $\mathrm{softmax}(z/T + b)$ with $b_0 = 0$; its MAP
version with 50 labels (\emph{MAP BCTS}) places Gaussian priors on $\log T$ and on each free bias, and its
\emph{MAP bias} variant keeps the source temperature and fits only $b$. The prior variances are the mean squared
oracle deviations measured on the other task, so nothing is fitted on the target being evaluated. A
\emph{prior adjustment} multiplies the source-TS probabilities by $\pi_Q / \pi_P$ and renormalises
\citep{saerens2002adjusting}, with $\pi_Q$ either estimated from the $n$ labels or, as a reference, the true target
calibration prior. Batch calibration \citep{zhou2024batch} needs no labels: it divides the source-TS probabilities by
their arithmetic mean over a batch of unlabelled target inputs and renormalises, which pulls the predicted class
proportions towards uniform. It therefore assumes a balanced target and acts on a batch rather than a single item.
Finally, a \emph{tuned threshold} adds the bias offset that maximises macro-F1 on the $n$ labelled items.

For every $T > 0$, $\mathrm{softmax}(z/T)$ preserves the order of the classes within an item, and for $K = 2$ it
preserves the order of the scores across items, since the logit difference is divided by the same positive constant.
Temperature scaling therefore changes no predicted class, and for $K = 2$ neither the AUROC nor the set of decisions
reachable by thresholding the score: a threshold calibrated on the scores, as in conformal risk control, flags the
same items. This is a fact about monotone maps, not an empirical finding. A bias term moves the boundary between
classes and changes decisions; for $K = 2$ it leaves the ranking of items intact.

\subsection{When a bias is the right repair}
\label{sec:background-location}

Consider the binary case, with $Z = z_1(X) - z_0(X)$ the source-TS log-odds and $\sigma(u) = 1/(1 + e^{-u})$. Two
assumptions describe the simplest shift that keeps the ranking intact. (A1) The score is calibrated on the source:
$P_P(Y = 1 \mid Z = u) = \sigma(u)$. (A2) Both classes move by the same amount: there is one $\delta \in \mathbb{R}$
such that, for each class $y$, the law of $Z$ given $Y = y$ under $Q$ is the law of $Z + \delta$ given $Y = y$ under
$P$; the class prior may move from $\pi_P$ to $\pi_Q$. The class-conditional laws are assumed to have densities
with respect to a common measure. Let $\lambda = \mathrm{logit}\, \pi_Q - \mathrm{logit}\, \pi_P$
and $\beta = \lambda - \delta$.

\begin{proposition}[A common location shift needs a bias]
\label{prop:location}
Under (A1) and (A2): (a) the target posterior is $P_Q(Y = 1 \mid Z = u) = \sigma(u + \beta)$, a bias at temperature 1;
(b) the AUROC of $Z$ is the same under $Q$ as under $P$; (c) every temperature map $\sigma(u/T)$, $T > 0$, keeps the
source decisions $\mathbb{1}[u > 0]$ while the Bayes rule is $\mathbb{1}[u > -\beta]$, and if $\beta \ne 0$, $Z$ has
no atoms and $\mathbb{E}_Q|Z| < \infty$, the target log-loss of the best temperature exceeds that of the bias map by a
positive amount; (d) with $K$ classes and a common shift $d \in \mathbb{R}^K$ of the logit vector, the target posterior
is $\mathrm{softmax}(s + b)$ with $b_k = \log(\pi_{Q,k}/\pi_{P,k}) - d_k$.
\end{proposition}

None of this is new as mathematics: parts (a) and (d) are the prior correction of \citet{saerens2002adjusting} with
a score shift added, and the link between a class prior and a decision threshold goes back at least to
\citet{elkan2001foundations}. The proof (Appendix~\ref{sec:appendix-proofs}) is a direct application of Bayes' rule.
The proposition is used here to name the assumption under which a bias, not a temperature, is the right repair, and to say what a fitted temperature means
when the assumption holds: a temperature fitted on target labels under (A2) is not 1, because it rises to hedge the
misplaced threshold. In a Gaussian model with $|\beta| = 2$ the oracle log temperature is 0.71--1.12, depending on
the class separation and the direction of the shift (Appendix~\ref{sec:appendix-simulation}). A measured
temperature shift can therefore be a bias shift read through the wrong parameter. The decomposition $\beta = \lambda - \delta$ also separates a change of the class
prior ($\lambda$) from a shift of the score ($\delta$). Label-shift estimators assume $\delta = 0$ and attribute every
change of the predicted proportions to the prior; when the score moves they estimate the wrong prior. Proposition~\ref{prop:location} is
an idealised limit: under (A2) the AUROC is retained exactly, whereas real threshold shifts lose some ranking. If the
classes move by different amounts, the target log-odds stay linear in $u$ in an equal-variance Gaussian model but with
another slope, so both a temperature and a bias are needed. On a labelled target sample, the slope that BCTS estimates
is the diagnostic: under (A2) its change of $\log T$ is close to zero.

\subsection{Quantities that describe a shift}
\label{sec:background-coordinates}

For each model and target we record, after source TS:
(i) the ranking \emph{AUROC} (macro one-vs-rest for $K = 3$) and its \emph{retention}, the ratio to the in-domain AUROC;
(ii) the \emph{predicted-marginal shift}, the total variation between the distribution of predicted labels on the
target inputs and on the source calibration split, which needs no target labels, and for comparison the total
variation between the target and source class priors, which does;
(iii) the \emph{oracle temperature change} $\log(T^{*}_{Q}/T_{P})$, where $T^{*}_{Q}$ maximises the likelihood of all
target calibration labels, flagged when the fit reaches the bound $T = 20$, that is, when the optimum is unbounded;
(iv) the \emph{BCTS coordinates}, the change of $\log T$ and the bias of BCTS fitted on all target calibration labels,
and the shares of the BCTS log-loss gain that the temperature alone and the bias alone recover; and
(v) the \emph{explained label entropy} $1 - \mathrm{NLL}_{Q}(\text{target TS}) / H(\pi_Q)$, whose value is at most 0
when even the oracle temperature is no better than predicting the target prior $\pi_Q$ for every input.

\subsection{Descriptive regimes}
\label{sec:background-regimes}

The primary analysis uses these quantities directly, over all 21 pairs. To describe the pairs we also use five labels,
which were defined after the results had been seen and are not a deployment diagnostic, since three of the quantities
above need target labels:
R0, \emph{no correction} (ranking retained, temperature and predicted proportions stable);
R1, \emph{temperature shift} (a finite temperature change with stable predicted proportions);
R2, \emph{threshold shift} (the predicted proportions move, target TS beats a constant predictor and the oracle
temperature is finite in most runs);
R3, \emph{weak ranking} (AUROC above 0.6 but target TS no better than a constant in most runs); and
R4, \emph{collapse} (an unbounded oracle temperature in every run). The operational thresholds are listed with the
regime table in Appendix~\ref{sec:appendix-regime-sensitivity}, which also varies them. The practical conclusions of
this article are stated over all pairs or over the pairs with a predicted-marginal shift above 0.10, and do not depend
on the labels.
